% !PRINTED 48
\weberpage{48}
%<<CONT
wodurch die Gleichungen (16) leicht in die Form gebracht werden:
\weq{17}{\begin{aligned}
\sqrt{x_1} &= -\dfrac{\sqrt{x_2}\sqrt{(1-x_3)(1-\varkappa^2x_3)} + \sqrt{x_3}\sqrt{(1-x_2)(1-\varkappa^2x_2)}}{1-\varkappa^2x_2x_3}, \\
\sqrt{1-x_1} &= \dfrac{\sqrt{1-x_2}\sqrt{1-x_3} - \sqrt{x_2x_3}\sqrt{(1-\varkappa^2x_2)(1-\varkappa^2x_3)}}{1-\varkappa^2x_2x_3}, \\
\sqrt{1-\varkappa^2x_1} &= \dfrac{\sqrt{1-\varkappa^2x_2}\sqrt{1-\varkappa^2x_3} - \varkappa^2\sqrt{x_2x_3}\sqrt{(1-x_2)(1-x_3)}}{1-\varkappa^2x_2x_3},
\end{aligned}}
und unsere Betrachtung lehrt, daß, wenn $x_1$ und die Wurzeln
$\sqrt{x_1}, \sqrt{1-x_2}, \sqrt{1-\varkappa^2x_1}$ % [sic] source lists \sqrt{1-x_2}; equations (17) and this x_1 list require \sqrt{1-x_1}.
durch diese Gleichungen als Funktionen
von $x_2, x_3$ bestimmt werden, die Differentialgleichung (1) identisch
befriedigt ist.

\begin{center}\rule{2cm}{0.4pt}\end{center}

Um auch für die Weierstrasssche Normalform das Additionstheorem
 in möglichst einfacher Gestalt zu erhalten, setze man
\[ f(x) = 4x^3 - g_2x - g_3, \]
und lasse in der Gleichung (11) (nach Division mit $x_4^2$) $x_4$ unendlich
 groß werden. Es ergibt sich alsdann das Integral der
Differentialgleichung (1)
\[ \dfrac{dx_1}{\sqrt{f(x_1)}} + \dfrac{dx_2}{\sqrt{f(x_2)}} + \dfrac{dx_3}{\sqrt{f(x_3)}} = 0 \]
in der Form
\weq{18}{\begin{vmatrix} 1, & x_1, & \sqrt{f(x_1)} \\ 1, & x_2, & \sqrt{f(x_2)} \\ 1, & x_3, & \sqrt{f(x_3)} \end{vmatrix} = 0,}
oder
\weq{19}{a + bx_1 + c\sqrt{f(x_1)} = 0,}
wenn
\weq{20}{\begin{aligned}
a &= x_2\sqrt{f(x_3)} - x_3\sqrt{f(x_2)} \\
b &= \sqrt{f(x_2)} - \sqrt{f(x_3)} \\
c &= x_3 - x_2.
\end{aligned}}
Es werden dann $x_1, x_2, x_3$ die Wurzeln der kubischen
Gleichung
\[ (a+bx)^2 - c^2f(x) = 0, \]
und wenn man hierin den Koeffizienten von $x^2$, geteilt durch
den von $x^3$ gleich der negativen Summe der Wurzeln setzt, so
erhält man
%CONT>>
